% figures/zipf.tex -- 词频—排名分布（双对数），Zipf 定律
%
% 数据：reports/data/csv/tasks/wordcount-default.csv（本轮 mse count，默认管线）
% 由 figures/data/make_figures_data.py 生成 figures/data/zipf.dat 与 zipf-meta.tex。
% 主文档只写 % figures/zipf.tex -- 词频—排名分布（双对数），Zipf 定律
%
% 数据：reports/data/csv/tasks/wordcount-default.csv（本轮 mse count，默认管线）
% 由 figures/data/make_figures_data.py 生成 figures/data/zipf.dat 与 zipf-meta.tex。
% 主文档只写 % figures/zipf.tex -- 词频—排名分布（双对数），Zipf 定律
%
% 数据：reports/data/csv/tasks/wordcount-default.csv（本轮 mse count，默认管线）
% 由 figures/data/make_figures_data.py 生成 figures/data/zipf.dat 与 zipf-meta.tex。
% 主文档只写 % figures/zipf.tex -- 词频—排名分布（双对数），Zipf 定律
%
% 数据：reports/data/csv/tasks/wordcount-default.csv（本轮 mse count，默认管线）
% 由 figures/data/make_figures_data.py 生成 figures/data/zipf.dat 与 zipf-meta.tex。
% 主文档只写 \input{figures/zipf}。
\input{figures/data/zipf-meta}
\begin{figure}[!htbp]
\centering
\begin{tikzpicture}
\begin{axis}[
  width=0.94\linewidth, height=6.0cm,
  xmode=log, ymode=log,
  xmin=1, xmax=40000, ymin=1, ymax=1e5,
  xlabel={rank (log scale)}, ylabel={term frequency tf (log scale)},
  xtick={1,10,100,1000,10000},
  ytick={1,10,100,1000,10000},
  tick label style={font=\scriptsize},
  label style={font=\small},
  grid=major, grid style={black!8},
  legend style={draw=none, fill=none, font=\scriptsize,
                at={(0.98,0.98)}, anchor=north east},
  legend cell align=left,
]
  % 实测 tf(rank)
  \addplot[only marks, mark=*, mark size=1.0pt, color=mseblue]
    table[x=rank, y=tf] {figures/data/zipf.dat};
  \addlegendentry{measured tf (\ZipfTerms{} terms)}
  % 1/rank 参考线（按 rank 1 处缩放对齐）
  \addplot[dashed, line width=0.8pt, color=msegray]
    table[x=rank, y=ref] {figures/data/zipf.dat};
  \addlegendentry{$1/\mathrm{rank}$ reference (scaled at rank 1)}
  % 最频词标注
  \node[font=\scriptsize, text=msegray, anchor=west] (topword)
    at (axis cs:1.15,42000) {the};
  \draw[->, draw=msegray, line width=0.4pt] (topword.west) -- (axis cs:1.0,3.1e4);
  % hapax 尾标注
  \node[font=\scriptsize, text=msegray, align=left, anchor=south west]
    at (rel axis cs:0.02,0.03)
    {hapax tail: \ZipfHapax{} terms (\ZipfHapaxPct\% of vocabulary)\\
     occur exactly once (tf $=1$)};
\end{axis}
\end{tikzpicture}
\caption{莎士比亚语料的词频—排名分布（双对数）：实测点近似一条直线，符合
Zipf 定律；\ZipfTerms{} 个词中有 \ZipfHapax{} 个只出现一次（占词表
\ZipfHapaxPct\%），构成需要倒排索引兜住的长尾。}%
\label{fig:zipf}
\end{figure}
。
% zipf-meta.tex -- generated by figures/data/make_figures_data.py
\newcommand{\ZipfTerms}{25\,705}
\newcommand{\ZipfHapax}{9\,753}
\newcommand{\ZipfHapaxPct}{37.9}

\begin{figure}[!htbp]
\centering
\begin{tikzpicture}
\begin{axis}[
  width=0.94\linewidth, height=6.0cm,
  xmode=log, ymode=log,
  xmin=1, xmax=40000, ymin=1, ymax=1e5,
  xlabel={rank (log scale)}, ylabel={term frequency tf (log scale)},
  xtick={1,10,100,1000,10000},
  ytick={1,10,100,1000,10000},
  tick label style={font=\scriptsize},
  label style={font=\small},
  grid=major, grid style={black!8},
  legend style={draw=none, fill=none, font=\scriptsize,
                at={(0.98,0.98)}, anchor=north east},
  legend cell align=left,
]
  % 实测 tf(rank)
  \addplot[only marks, mark=*, mark size=1.0pt, color=mseblue]
    table[x=rank, y=tf] {figures/data/zipf.dat};
  \addlegendentry{measured tf (\ZipfTerms{} terms)}
  % 1/rank 参考线（按 rank 1 处缩放对齐）
  \addplot[dashed, line width=0.8pt, color=msegray]
    table[x=rank, y=ref] {figures/data/zipf.dat};
  \addlegendentry{$1/\mathrm{rank}$ reference (scaled at rank 1)}
  % 最频词标注
  \node[font=\scriptsize, text=msegray, anchor=west] (topword)
    at (axis cs:1.15,42000) {the};
  \draw[->, draw=msegray, line width=0.4pt] (topword.west) -- (axis cs:1.0,3.1e4);
  % hapax 尾标注
  \node[font=\scriptsize, text=msegray, align=left, anchor=south west]
    at (rel axis cs:0.02,0.03)
    {hapax tail: \ZipfHapax{} terms (\ZipfHapaxPct\% of vocabulary)\\
     occur exactly once (tf $=1$)};
\end{axis}
\end{tikzpicture}
\caption{莎士比亚语料的词频—排名分布（双对数）：实测点近似一条直线，符合
Zipf 定律；\ZipfTerms{} 个词中有 \ZipfHapax{} 个只出现一次（占词表
\ZipfHapaxPct\%），构成需要倒排索引兜住的长尾。}%
\label{fig:zipf}
\end{figure}
。
% zipf-meta.tex -- generated by figures/data/make_figures_data.py
\newcommand{\ZipfTerms}{25\,705}
\newcommand{\ZipfHapax}{9\,753}
\newcommand{\ZipfHapaxPct}{37.9}

\begin{figure}[!htbp]
\centering
\begin{tikzpicture}
\begin{axis}[
  width=0.94\linewidth, height=6.0cm,
  xmode=log, ymode=log,
  xmin=1, xmax=40000, ymin=1, ymax=1e5,
  xlabel={rank (log scale)}, ylabel={term frequency tf (log scale)},
  xtick={1,10,100,1000,10000},
  ytick={1,10,100,1000,10000},
  tick label style={font=\scriptsize},
  label style={font=\small},
  grid=major, grid style={black!8},
  legend style={draw=none, fill=none, font=\scriptsize,
                at={(0.98,0.98)}, anchor=north east},
  legend cell align=left,
]
  % 实测 tf(rank)
  \addplot[only marks, mark=*, mark size=1.0pt, color=mseblue]
    table[x=rank, y=tf] {figures/data/zipf.dat};
  \addlegendentry{measured tf (\ZipfTerms{} terms)}
  % 1/rank 参考线（按 rank 1 处缩放对齐）
  \addplot[dashed, line width=0.8pt, color=msegray]
    table[x=rank, y=ref] {figures/data/zipf.dat};
  \addlegendentry{$1/\mathrm{rank}$ reference (scaled at rank 1)}
  % 最频词标注
  \node[font=\scriptsize, text=msegray, anchor=west] (topword)
    at (axis cs:1.15,42000) {the};
  \draw[->, draw=msegray, line width=0.4pt] (topword.west) -- (axis cs:1.0,3.1e4);
  % hapax 尾标注
  \node[font=\scriptsize, text=msegray, align=left, anchor=south west]
    at (rel axis cs:0.02,0.03)
    {hapax tail: \ZipfHapax{} terms (\ZipfHapaxPct\% of vocabulary)\\
     occur exactly once (tf $=1$)};
\end{axis}
\end{tikzpicture}
\caption{莎士比亚语料的词频—排名分布（双对数）：实测点近似一条直线，符合
Zipf 定律；\ZipfTerms{} 个词中有 \ZipfHapax{} 个只出现一次（占词表
\ZipfHapaxPct\%），构成需要倒排索引兜住的长尾。}%
\label{fig:zipf}
\end{figure}
。
% zipf-meta.tex -- generated by figures/data/make_figures_data.py
\newcommand{\ZipfTerms}{25\,705}
\newcommand{\ZipfHapax}{9\,753}
\newcommand{\ZipfHapaxPct}{37.9}

\begin{figure}[!htbp]
\centering
\begin{tikzpicture}
\begin{axis}[
  width=0.94\linewidth, height=6.0cm,
  xmode=log, ymode=log,
  xmin=1, xmax=40000, ymin=1, ymax=1e5,
  xlabel={rank (log scale)}, ylabel={term frequency tf (log scale)},
  xtick={1,10,100,1000,10000},
  ytick={1,10,100,1000,10000},
  tick label style={font=\scriptsize},
  label style={font=\small},
  grid=major, grid style={black!8},
  legend style={draw=none, fill=none, font=\scriptsize,
                at={(0.98,0.98)}, anchor=north east},
  legend cell align=left,
]
  % 实测 tf(rank)
  \addplot[only marks, mark=*, mark size=1.0pt, color=mseblue]
    table[x=rank, y=tf] {figures/data/zipf.dat};
  \addlegendentry{measured tf (\ZipfTerms{} terms)}
  % 1/rank 参考线（按 rank 1 处缩放对齐）
  \addplot[dashed, line width=0.8pt, color=msegray]
    table[x=rank, y=ref] {figures/data/zipf.dat};
  \addlegendentry{$1/\mathrm{rank}$ reference (scaled at rank 1)}
  % 最频词标注
  \node[font=\scriptsize, text=msegray, anchor=west] (topword)
    at (axis cs:1.15,42000) {the};
  \draw[->, draw=msegray, line width=0.4pt] (topword.west) -- (axis cs:1.0,3.1e4);
  % hapax 尾标注
  \node[font=\scriptsize, text=msegray, align=left, anchor=south west]
    at (rel axis cs:0.02,0.03)
    {hapax tail: \ZipfHapax{} terms (\ZipfHapaxPct\% of vocabulary)\\
     occur exactly once (tf $=1$)};
\end{axis}
\end{tikzpicture}
\caption{莎士比亚语料的词频—排名分布（双对数）：实测点近似一条直线，符合
Zipf 定律；\ZipfTerms{} 个词中有 \ZipfHapax{} 个只出现一次（占词表
\ZipfHapaxPct\%），构成需要倒排索引兜住的长尾。}%
\label{fig:zipf}
\end{figure}
